\section*{Bab \ref{ch:b3:computational-chemistry} --- Kimia Komputasi: Hartree--Fock dan DFT}

\begin{solution}{exo:b3:computational-chemistry:1}
$\dd E/\dd\alpha = \frac32 - \sqrt{2/\pi\alpha} = 0$ memberikan $\alpha = 8/9\pi = 0.283\,
a_0^{-2}$ dan $E = \frac32\alpha - 2\sqrt{2\alpha/\pi} = -4/3\pi = \qty{-0.4244}{\hartree}$,
15\,\% di atas nilai eksak $-0.5$: satu fungsi Gauss tidak mempunyai titik
runcing maupun ekor fungsi $1s$.
\end{solution}

\begin{solution}{exo:b3:computational-chemistry:2}
STO-3G: 5 fungsi per C ($1s$, $2s$, tiga $2p$) dan 1 per H: $6 \times 5 + 6 =
36$. 6-31G(d): per C, 1 (inti dalam) + $2 \times 4$ (valensi terpisah) + 6 ($d$)
$= 15$; per H, 2: $6 \times 15 + 6 \times 2 = 102$.
\end{solution}

\begin{solution}{exo:b3:computational-chemistry:3}
Hamiltonian elektronik tidak memuat massa inti: dalam \omterm{def:b3:computational-chemistry:born-oppenheimer}{pendekatan
Born--Oppenheimer}, PES-nya, sehingga juga $r_e$ dan $k = U''(r_e)$, sama.
$\tilde\omega = \sqrt{k/\mu}/2\pi c$ bergantung pada \omterm{def:b3:quantum-model-systems:oscillator}{massa tereduksi}, yang
menjadi dua kali lipat: rasionya $\sqrt2 = 1.414$ (terukur 1.413).
\end{solution}

\begin{solution}{exo:b3:computational-chemistry:4}
Delapan atom memberikan $3N - 6 = 18$ vibrasi: 17 real dan satu imajiner.
Satu \omterm{def:b3:quantum-model-systems:operator}{nilai eigen} Hessian negatif: titik pelana orde pertama, sebuah struktur
transisi (di sini, misalnya, dari sebuah rotasi atau sebuah eliminasi).
\end{solution}

\begin{solution}{exo:b3:computational-chemistry:5}
$\zeta = 27/16$, $E = -(27/16)^2 = \qty{-2.8477}{\hartree} = \qty{-77.49}{eV}$.
Galatnya $-2.8477 + 2.9034 = \qty{0.0557}{\hartree} = \qty{1.52}{eV}$, yaitu
1.9\,\% dari energi total --- tetapi lebih besar daripada banyak energi reaksi.
\end{solution}

\begin{solution}{exo:b3:computational-chemistry:6}
$(-1 - E)(-0.5 - E) - (-0.2 - 0.3E)^2 = 0$, yaitu $0.91E^2 + 1.38E + 0.46 = 0$:
$E = \qty{-1.0217}{\hartree}$ dan \qty{-0.4947}{\hartree}. Ya: mencampurkan
fungsi kedua menurunkan energi di bawah $H_{11}$, sebagaimana diizinkan \omterm{thm:b3:computational-chemistry:variational}{prinsip
variasi}.
\end{solution}

\begin{solution}{exo:b3:computational-chemistry:7}
$(102/36)^4 = 64$: sekitar 640 menit, kira-kira 11 jam.
\end{solution}

\begin{solution}{exo:b3:computational-chemistry:8}
$0.5782 \times 27.211 = \qty{15.73}{eV}$, 0.30 eV di atas nilai terukur.
Orbital yang dibekukan: ion yang sebenarnya berelaksasi dan energinya lebih
rendah, sehingga nilai Koopmans terlalu tinggi. Korelasi yang hilang: molekul
netral (dua elektron) mempunyai \omterm{def:b3:computational-chemistry:correlation}{energi korelasi} yang lebih besar daripada ion
(satu elektron), yang membuat nilai sebenarnya lebih besar. Di sini galat
yang pertama lebih menentukan.
\end{solution}

\begin{solution}{exo:b3:computational-chemistry:9}
$-0.5960 - 2(-0.4666) = \qty{0.3372}{\hartree} = \qty{9.18}{eV}$ di atas dua atom.
Fungsi RHF mempertahankan 50\,\% \ce{H+ H-}, dan memisahkan sebuah proton dan
sebuah hidrida memerlukan selisih antara energi ionisasi H dan afinitas
elektron H, yang ikut terambil dalam rata-rata energinya.
\end{solution}

\begin{solution}{exo:b3:computational-chemistry:10}
$E(c_1^2 + 2c_1c_2S + c_2^2) = c_1^2H_{11} + 2c_1c_2H_{12} + c_2^2H_{22}$.
Dengan menurunkan terhadap $c_1$ dan $\partial E/\partial c_1 = 0$:
$E(2c_1 + 2c_2S) = 2c_1H_{11} + 2c_2H_{12}$, yaitu $(H_{11} - E)c_1 + (H_{12} -
ES)c_2 = 0$; dengan cara yang sama $(H_{12} - ES)c_1 + (H_{22} - E)c_2 = 0$.
\end{solution}

\begin{solution}{exo:b3:computational-chemistry:11}
$k \approx [E(1.30) - 2E(1.35) + E(1.40)]/(0.05)^2 = 0.001417/0.0025 =
\qty{0.567}{\hartree\per\bohr\squared} = \qty{882}{N/m}$. Dengan $\mu = m_H/2 =
\qty{8.37e-28}{kg}$, $\omega = \sqrt{k/\mu} = \qty{1.03e15}{s^{-1}}$ dan
$\tilde\omega = \omega/2\pi c = \qty{5452}{cm^{-1}}$, 24\,\% di atas
$\tilde\omega_e$ yang terukur: kurva RHF basis minimal terlalu curam (tanpa
korelasi, basis terlalu kecil), itulah sebabnya frekuensi hasil perhitungan
diskalakan turun.
\end{solution}

\begin{solution}{exo:b3:computational-chemistry:12}
Selisih beberapa \unit{kJ/mol} yang melibatkan ikatan hidrogen memerlukan
korelasi dan dispersi: fungsional hibrida dengan koreksi dispersi, atau lebih
baik metode fungsi gelombang berkorelasi untuk energi akhirnya, dengan basis
terpolarisasi bervalensi terpisah rangkap tiga yang memuat \omterm{def:b3:computational-chemistry:split-valence}{fungsi difus}.
Optimasikan kedua konformer pada tingkat yang sama, pastikan keduanya
minimum melalui frekuensinya (semua real), tambahkan \omterm{def:b3:quantum-model-systems:oscillator}{energi titik nol}, uji
basisnya dengan satu perhitungan yang lebih besar, dan bandingkan dengan
sistem sejenis yang energi konformasinya diketahui.
\end{solution}

\begin{solution}{pb:b3:computational-chemistry:1}
\textbf{1.} Setiap \omterm{def:b3:computational-chemistry:basis}{orbital tipe Slater} dalam basis minimal diganti dengan
kontraksi tetap tiga fungsi Gauss yang dicocokkan dengannya.
\textbf{2.} Inti helium menarik elektron-elektronnya lebih kuat, sehingga
fungsi $1s$-nya lebih kompak (eksponen lebih besar).
$6.3624/3.4253 = 1.857 = (1.69/1.24)^2$.
\textbf{3.} Dua fungsi basis, dua elektron, satu orbital terisi (dan satu
orbital virtual).
\textbf{4.} $V_{\mathrm{nn}} = Z_{\mathrm{He}}Z_{\mathrm H}/R = 2/1.4632 =
\qty{1.3669}{\hartree}$.
\textbf{5.} Kedua fungsi itu bertumpang-tindih kuat (54\,\%): atom-atomnya
cukup dekat untuk berikatan.
\textbf{6.} $h_{11}$ adalah energi sebuah elektron dalam fungsi He dengan
kedua inti tetapi tanpa elektron lain; inti He (bermuatan 2) mengikatnya jauh
lebih erat.
\textbf{7.} $(h_{11} - \varepsilon)(h_{22} - \varepsilon) - (h_{12} - \varepsilon S)^2 = 0$:
$0.71185\varepsilon^2 + 2.79292\varepsilon + 2.44969 = 0$.
\textbf{8.} $\varepsilon = \qty{-2.600}{\hartree}$ dan \qty{-1.324}{\hartree}.
\textbf{9.} $\mathbf h$ mengabaikan tolakan antara kedua elektron; orbital
tebakan inti terlalu mengerut dan terlalu rendah.
\textbf{10.} $F_{\mu\nu} = h_{\mu\nu} + \sum_{\lambda\sigma}P_{\lambda\sigma}[(\mu\nu|\sigma\lambda) -
\frac12(\mu\lambda|\sigma\nu)]$, dengan $P_{\lambda\sigma} = 2C_{\lambda1}C_{\sigma1}$.
\textbf{11.} Tolakan Coulomb setiap elektron oleh awan muatan elektron yang
lain (dan, secara umum, pertukaran).
\textbf{12.} $\mathbf F$ bergantung pada $\mathbf P$, yang bergantung pada
orbital-orbital yang diperoleh dari $\mathbf F$: persamaannya taklinear dan
diselesaikan dengan pendekatan berturut-turut sampai swakonsisten.
\textbf{13.} $c_1^2 + c_2^2 + 2c_1c_2S = 0.7684 + 0.0410 + 0.1906 = 1.0000$.
\textbf{14.} He: $1.5369 + 0.1906 = 1.727$; H: $0.0820 + 0.1906 = 0.273$
elektron.
\textbf{15.} Muatan: He $+0.27$, H $+0.73$. Muatan positifnya sebagian besar
berada pada hidrogen: \ce{HeH+} adalah proton yang terikat pada atom helium,
\ce{He + H+}, sesuai harapan karena energi ionisasi He (24.6 eV) jauh melebihi
energi ionisasi H (13.6 eV).
\textbf{16.} $-\varepsilon_1 = \qty{1.633}{\hartree} = \qty{44.4}{eV}$.
\textbf{17.} $E_{\mathrm{el}} = -2.8418 - 1.3669 = \qty{-4.2087}{\hartree}$.
\textbf{18.} Tiga siklus mencapai \qty{-2.8418}{\hartree}. Setiap siklus
memberikan sebuah determinan, yang energinya merupakan batas atas bagi energi
Hartree--Fock yang konvergen (\omterm{thm:b3:computational-chemistry:variational}{prinsip variasi}), dan iterasi-iterasinya
memperbaikinya.
\textbf{19.} Basis dengan $\zeta = 2.0925$ memberikan energi yang lebih rendah,
sehingga basis itu lebih baik untuk molekul ini (\omterm{thm:b3:computational-chemistry:variational}{prinsip variasi}): fungsi
helium lebih kompak dalam \ce{HeH+} daripada dalam atom bebas, yang nilai
$\zeta = 1.69$-nya ditiru oleh basis standar.
\textbf{20.} Orbital antiikatan kosong $\sigma^*$; energinya akan mendekati
minus afinitas elektron \ce{HeH+} (dengan buruk, dalam basis sekecil ini).
\textbf{21.} Nol (tanpa elektron); $\qty{-2.8078}{\hartree}$.
\textbf{22.} $-2.8078 - (-2.8418) = \qty{0.0340}{\hartree} = \qty{89}{kJ/mol}$.
\textbf{23.} Sekitar setengah dari nilai terukur: basis minimal tidak dapat
memolarisasi kerapatan helium ke arah proton (tidak ada fungsi $p$), sehingga
ikatannya terlalu lemah. (Koreksi titik nol dan koreksi \qty{298}{K} kecil
dibandingkan dengan ini.)
\textbf{24.} $-(24.5874 + 54.4178)/27.211 = \qty{-2.9034}{\hartree}$; atom STO-3G
\qty{0.0956}{\hartree} (\qty{2.6}{eV}) terlalu tinggi.
\textbf{25.} Geometri bergantung pada cara energi berubah terhadap $R$;
sebagian besar galatnya, yang terpusat pada kulit dalam, hampir sama pada
setiap geometri dan saling meniadakan.
\textbf{26.} \textbf{$E(\ce{HeH+}) = \qty{-2.8418}{\hartree}$ pada RHF/STO-3G
(basis standar), \qty{1.4632}{\bohr}} (dan \qty{-2.8607}{\hartree} dengan nilai
klasik $\zeta_{\mathrm{He}} = 2.0925$).
\end{solution}
